\documentclass[../../main.tex]{subfiles}
\begin{document}

\begin{center}
  {\fontsize{48pt}{58pt}\selectfont \textbf{1993}\normalsize\textbf{年度}}
\end{center}

\begin{table}[htb]
  \centering
  \caption{出題テーマと難易度}
  \begin{tabular}{|c|p{6.2cm}|c|c|c|}
    \hline
    \textbf{問題番号} & \textbf{テーマ} & \textbf{思考量} & \textbf{計算量} & \textbf{総合} \\
    \hline
    第1問 & 三次関数の係数と面積条件 & 普 & 難 & 普 \\
    \hline
    第2問 & 三次方程式の実根の個数 & 普 & 普 & 普 \\
    \hline
    第3問 & 円周上の4点から作る重心の列 & 難 & 難 & 難 \\
    \hline
    第4問 & 指数関数の不等式 & 難 & 普 & 難 \\
    \hline
    第5問 & 重複を許す数字の並べ方 & 普 & 普 & 普 \\
    \hline
    第6問 & 二つの分数関数の大小関係 & 難 & 難 & 難 \\
    \hline
  \end{tabular}
\end{table}

\begin{table}[htb]
  \centering
  \caption{解答}
  \renewcommand{\arraystretch}{1.4}
  \begin{tabular}{|c|c|p{9cm}|}
    \hline
    \textbf{問題番号} & \textbf{小問} & \textbf{答え} \\
    \hline
    第1問 & -- & $(a,b,c)=(3,-7,3)$ \\
    \hline
    第2問 & (1) & $a>0$かつ$-2a\sqrt{a}<t<2a\sqrt{a}$ \\
    \cline{2-3}
    & (2) & $a>\dfrac{\sqrt{3}}{2}$ \\
    \hline
    第3問 & (1) & $\overrightarrow{OP_1}=\dfrac{\vec{a}+\vec{b}+\vec{c}+\vec{d}}{3}$ \\
    \cline{2-3}
    & (2) & $\overrightarrow{P_nP_{n+1}}=-\left(-\dfrac{1}{3}\right)^{n+1}(\vec{a}+\vec{b}+\vec{c}+\vec{d})$ \\
    \cline{2-3}
    & (3) & $\overrightarrow{OQ}=\dfrac{\vec{a}+\vec{b}+\vec{c}+\vec{d}}{4}$ \\
    \hline
    第4問 & -- & $a=e$ \\
    \hline
    第5問 & -- & $n(15n^2-20n+6)$通り \\
    \hline
    第6問 & -- & $a_1>a_2$かつ，$c_1=c_2$なら$b_2-a_2c_2=b_1-a_1c_1$，$c_1\ne c_2$なら$b_1=a_1c_1$，$b_2=a_2c_2$ \\
    \hline
  \end{tabular}
\end{table}

\paragraph{試験形態}
6問．試験時間は要確認．

\end{document}
